% !TEX program = xelatex
% Báo cáo Assignment 1 - BrewLite. Biên dịch: latexmk -xelatex BaoCao_Assignment1.tex
\documentclass[12pt, a4paper]{article}
\newcommand{\mssvLong}{MSSV: ..........}   % TODO: điền MSSV
\newcommand{\mssvTai}{MSSV: ..........}    % TODO: điền MSSV
\newcommand{\sgufooter}{Bài tập lớn môn Công nghệ phần mềm -- Assignment 1 -- HK1 2026-2027}
% =====================================================================
%  Khung báo cáo theo mẫu SGU (Assignment_report.tex), chuyển sang XeLaTeX.
%  Dùng: % =====================================================================
%  Khung báo cáo theo mẫu SGU (Assignment_report.tex), chuyển sang XeLaTeX.
%  Dùng: % =====================================================================
%  Khung báo cáo theo mẫu SGU (Assignment_report.tex), chuyển sang XeLaTeX.
%  Dùng: \input{sgu} trong phần preamble, rồi gọi \sgutitle{...} sau \begin{document}.
%  Biên dịch: latexmk -xelatex <tên>.tex
% =====================================================================
\usepackage[a4paper, top=2.6cm, bottom=2.4cm, left=2.5cm, right=2.2cm,
            headheight=40pt, headsep=0.6cm]{geometry}
\usepackage{fontspec}
\usepackage[vietnamese, provide=*]{babel}
\IfFontExistsTF{Times New Roman}{\setmainfont{Times New Roman}}{\setmainfont{TeX Gyre Termes}}
\setsansfont{TeX Gyre Heros}
\IfFontExistsTF{Liberation Mono}{\setmonofont{Liberation Mono}[Scale=0.88]}{\setmonofont{FreeMono}}
\usepackage{microtype}
\usepackage{setspace}
\setstretch{1.15}

\usepackage{amsmath, amssymb}
\usepackage[table]{xcolor}
\usepackage{array, tabularx, booktabs, colortbl, multirow}
\usepackage{enumitem}
\usepackage{graphicx}
\usepackage{float}
\usepackage{caption}
\usepackage{fancyhdr}
\usepackage{lastpage}
\usepackage{tikz}
\usetikzlibrary{arrows.meta, positioning, shapes.geometric, fit, calc}
\usepackage[most]{tcolorbox}
\usepackage{fvextra}
\usepackage{xurl}
\usepackage[colorlinks=true, linkcolor=black, urlcolor=blue, citecolor=black,
            bookmarksnumbered=true]{hyperref}

\definecolor{sgugray}{HTML}{F3F5F8}
\definecolor{sguline}{HTML}{C9D1DC}
\definecolor{dkgreen}{rgb}{0,0.5,0}
\definecolor{mauve}{rgb}{0.58,0,0.82}

\setlength{\parindent}{0pt}
\setlength{\parskip}{5pt plus 1pt}
\setlength{\emergencystretch}{3em}
\setcounter{secnumdepth}{3}
\setcounter{tocdepth}{3}
\renewcommand{\arraystretch}{1.25}
\setlist{itemsep=2pt, topsep=3pt, parsep=0pt}
\captionsetup{font=small, labelfont=bf}

\usepackage{titlesec}
\titleformat{\section}{\sffamily\Large\bfseries}{\thesection}{0.6em}{}
\titleformat{\subsection}{\sffamily\large\bfseries}{\thesubsection}{0.6em}{}
\titleformat{\subsubsection}{\sffamily\normalsize\bfseries}{\thesubsubsection}{0.6em}{}

% Header / footer đúng mẫu SGU
\pagestyle{fancy}
\fancyhf{}
\fancyhead[L]{%
  \begin{tabular}{@{}rl@{}}
    \includegraphics[height=8mm]{logoITSGUsmall.png}\hspace{4pt} &
    \begin{tabular}{@{}l@{}}
      {\ttfamily\bfseries\small Trường Đại học Sài Gòn}\\
      {\ttfamily\bfseries\small Khoa Công Nghệ Thông Tin}
    \end{tabular}
  \end{tabular}}
\fancyfoot[L]{\scriptsize\ttfamily \sgufooter}
\fancyfoot[R]{\scriptsize\ttfamily Trang \thepage/\pageref{LastPage}}
\renewcommand{\headrulewidth}{0.3pt}
\renewcommand{\footrulewidth}{0.3pt}

\DefineVerbatimEnvironment{code}{Verbatim}{fontsize=\footnotesize, breaklines=true, frame=lines, framerule=0.3pt, rulecolor=\color{sguline}, numbers=left, numbersep=6pt, xleftmargin=14pt}
\newcommand{\codefile}[3]{\VerbatimInput[fontsize=\footnotesize, breaklines=true, frame=lines, framerule=0.3pt, rulecolor=\color{sguline}, numbers=left, numbersep=6pt, xleftmargin=14pt, firstline=#1, lastline=#2, firstnumber=#1]{#3}}

\newcommand{\fpath}[1]{\nolinkurl{#1}}
\newtcolorbox{ghichu}[1][]{enhanced, breakable, colback=sgugray, colframe=sguline,
  boxrule=0.5pt, arc=2pt, left=8pt, right=8pt, top=5pt, bottom=5pt,
  before skip=8pt, after skip=8pt, #1}

% Trang bìa theo mẫu SGU.  #1 đề tài, #2 loại bài, #3 tháng/năm
\newcommand{\sgutitle}[3]{%
\begin{titlepage}
\begin{center}
TRƯỜNG ĐẠI HỌC SÀI GÒN \\
KHOA CÔNG NGHỆ THÔNG TIN
\end{center}
\vspace{1cm}
\begin{center}\includegraphics[width=3cm]{logoITSGU.png}\end{center}
\vspace{1cm}
\begin{center}
\begin{tabular}{c}
  \multicolumn{1}{l}{\textbf{{\Large CÔNG NGHỆ PHẦN MỀM}}}\\
  ~~\\ \hline \\
  \multicolumn{1}{l}{\textbf{{\Large #2}}}\\ \\
  \textbf{{\Huge #1}}\\ \\ \hline
\end{tabular}
\end{center}
\vspace{2.5cm}
\begin{center}
\begin{tabular}{rrl}
  \hspace{4cm} & GVHD: & Từ Lãng Phiêu\\
  & SV: & \tenLong{} -- \mssvLong\\
  & & \tenTai{} -- \mssvTai\\
\end{tabular}
\end{center}
\vfill
\begin{center}{\footnotesize TP. HỒ CHÍ MINH, #3}\end{center}
\end{titlepage}
\thispagestyle{empty}
\newpage
\tableofcontents
\newpage}
 trong phần preamble, rồi gọi \sgutitle{...} sau \begin{document}.
%  Biên dịch: latexmk -xelatex <tên>.tex
% =====================================================================
\usepackage[a4paper, top=2.6cm, bottom=2.4cm, left=2.5cm, right=2.2cm,
            headheight=40pt, headsep=0.6cm]{geometry}
\usepackage{fontspec}
\usepackage[vietnamese, provide=*]{babel}
\IfFontExistsTF{Times New Roman}{\setmainfont{Times New Roman}}{\setmainfont{TeX Gyre Termes}}
\setsansfont{TeX Gyre Heros}
\IfFontExistsTF{Liberation Mono}{\setmonofont{Liberation Mono}[Scale=0.88]}{\setmonofont{FreeMono}}
\usepackage{microtype}
\usepackage{setspace}
\setstretch{1.15}

\usepackage{amsmath, amssymb}
\usepackage[table]{xcolor}
\usepackage{array, tabularx, booktabs, colortbl, multirow}
\usepackage{enumitem}
\usepackage{graphicx}
\usepackage{float}
\usepackage{caption}
\usepackage{fancyhdr}
\usepackage{lastpage}
\usepackage{tikz}
\usetikzlibrary{arrows.meta, positioning, shapes.geometric, fit, calc}
\usepackage[most]{tcolorbox}
\usepackage{fvextra}
\usepackage{xurl}
\usepackage[colorlinks=true, linkcolor=black, urlcolor=blue, citecolor=black,
            bookmarksnumbered=true]{hyperref}

\definecolor{sgugray}{HTML}{F3F5F8}
\definecolor{sguline}{HTML}{C9D1DC}
\definecolor{dkgreen}{rgb}{0,0.5,0}
\definecolor{mauve}{rgb}{0.58,0,0.82}

\setlength{\parindent}{0pt}
\setlength{\parskip}{5pt plus 1pt}
\setlength{\emergencystretch}{3em}
\setcounter{secnumdepth}{3}
\setcounter{tocdepth}{3}
\renewcommand{\arraystretch}{1.25}
\setlist{itemsep=2pt, topsep=3pt, parsep=0pt}
\captionsetup{font=small, labelfont=bf}

\usepackage{titlesec}
\titleformat{\section}{\sffamily\Large\bfseries}{\thesection}{0.6em}{}
\titleformat{\subsection}{\sffamily\large\bfseries}{\thesubsection}{0.6em}{}
\titleformat{\subsubsection}{\sffamily\normalsize\bfseries}{\thesubsubsection}{0.6em}{}

% Header / footer đúng mẫu SGU
\pagestyle{fancy}
\fancyhf{}
\fancyhead[L]{%
  \begin{tabular}{@{}rl@{}}
    \includegraphics[height=8mm]{logoITSGUsmall.png}\hspace{4pt} &
    \begin{tabular}{@{}l@{}}
      {\ttfamily\bfseries\small Trường Đại học Sài Gòn}\\
      {\ttfamily\bfseries\small Khoa Công Nghệ Thông Tin}
    \end{tabular}
  \end{tabular}}
\fancyfoot[L]{\scriptsize\ttfamily \sgufooter}
\fancyfoot[R]{\scriptsize\ttfamily Trang \thepage/\pageref{LastPage}}
\renewcommand{\headrulewidth}{0.3pt}
\renewcommand{\footrulewidth}{0.3pt}

\DefineVerbatimEnvironment{code}{Verbatim}{fontsize=\footnotesize, breaklines=true, frame=lines, framerule=0.3pt, rulecolor=\color{sguline}, numbers=left, numbersep=6pt, xleftmargin=14pt}
\newcommand{\codefile}[3]{\VerbatimInput[fontsize=\footnotesize, breaklines=true, frame=lines, framerule=0.3pt, rulecolor=\color{sguline}, numbers=left, numbersep=6pt, xleftmargin=14pt, firstline=#1, lastline=#2, firstnumber=#1]{#3}}

\newcommand{\fpath}[1]{\nolinkurl{#1}}
\newtcolorbox{ghichu}[1][]{enhanced, breakable, colback=sgugray, colframe=sguline,
  boxrule=0.5pt, arc=2pt, left=8pt, right=8pt, top=5pt, bottom=5pt,
  before skip=8pt, after skip=8pt, #1}

% Trang bìa theo mẫu SGU.  #1 đề tài, #2 loại bài, #3 tháng/năm
\newcommand{\sgutitle}[3]{%
\begin{titlepage}
\begin{center}
TRƯỜNG ĐẠI HỌC SÀI GÒN \\
KHOA CÔNG NGHỆ THÔNG TIN
\end{center}
\vspace{1cm}
\begin{center}\includegraphics[width=3cm]{logoITSGU.png}\end{center}
\vspace{1cm}
\begin{center}
\begin{tabular}{c}
  \multicolumn{1}{l}{\textbf{{\Large CÔNG NGHỆ PHẦN MỀM}}}\\
  ~~\\ \hline \\
  \multicolumn{1}{l}{\textbf{{\Large #2}}}\\ \\
  \textbf{{\Huge #1}}\\ \\ \hline
\end{tabular}
\end{center}
\vspace{2.5cm}
\begin{center}
\begin{tabular}{rrl}
  \hspace{4cm} & GVHD: & Từ Lãng Phiêu\\
  & SV: & \tenLong{} -- \mssvLong\\
  & & \tenTai{} -- \mssvTai\\
\end{tabular}
\end{center}
\vfill
\begin{center}{\footnotesize TP. HỒ CHÍ MINH, #3}\end{center}
\end{titlepage}
\thispagestyle{empty}
\newpage
\tableofcontents
\newpage}
 trong phần preamble, rồi gọi \sgutitle{...} sau \begin{document}.
%  Biên dịch: latexmk -xelatex <tên>.tex
% =====================================================================
\usepackage[a4paper, top=2.6cm, bottom=2.4cm, left=2.5cm, right=2.2cm,
            headheight=40pt, headsep=0.6cm]{geometry}
\usepackage{fontspec}
\usepackage[vietnamese, provide=*]{babel}
\IfFontExistsTF{Times New Roman}{\setmainfont{Times New Roman}}{\setmainfont{TeX Gyre Termes}}
\setsansfont{TeX Gyre Heros}
\IfFontExistsTF{Liberation Mono}{\setmonofont{Liberation Mono}[Scale=0.88]}{\setmonofont{FreeMono}}
\usepackage{microtype}
\usepackage{setspace}
\setstretch{1.15}

\usepackage{amsmath, amssymb}
\usepackage[table]{xcolor}
\usepackage{array, tabularx, booktabs, colortbl, multirow}
\usepackage{enumitem}
\usepackage{graphicx}
\usepackage{float}
\usepackage{caption}
\usepackage{fancyhdr}
\usepackage{lastpage}
\usepackage{tikz}
\usetikzlibrary{arrows.meta, positioning, shapes.geometric, fit, calc}
\usepackage[most]{tcolorbox}
\usepackage{fvextra}
\usepackage{xurl}
\usepackage[colorlinks=true, linkcolor=black, urlcolor=blue, citecolor=black,
            bookmarksnumbered=true]{hyperref}

\definecolor{sgugray}{HTML}{F3F5F8}
\definecolor{sguline}{HTML}{C9D1DC}
\definecolor{dkgreen}{rgb}{0,0.5,0}
\definecolor{mauve}{rgb}{0.58,0,0.82}

\setlength{\parindent}{0pt}
\setlength{\parskip}{5pt plus 1pt}
\setlength{\emergencystretch}{3em}
\setcounter{secnumdepth}{3}
\setcounter{tocdepth}{3}
\renewcommand{\arraystretch}{1.25}
\setlist{itemsep=2pt, topsep=3pt, parsep=0pt}
\captionsetup{font=small, labelfont=bf}

\usepackage{titlesec}
\titleformat{\section}{\sffamily\Large\bfseries}{\thesection}{0.6em}{}
\titleformat{\subsection}{\sffamily\large\bfseries}{\thesubsection}{0.6em}{}
\titleformat{\subsubsection}{\sffamily\normalsize\bfseries}{\thesubsubsection}{0.6em}{}

% Header / footer đúng mẫu SGU
\pagestyle{fancy}
\fancyhf{}
\fancyhead[L]{%
  \begin{tabular}{@{}rl@{}}
    \includegraphics[height=8mm]{logoITSGUsmall.png}\hspace{4pt} &
    \begin{tabular}{@{}l@{}}
      {\ttfamily\bfseries\small Trường Đại học Sài Gòn}\\
      {\ttfamily\bfseries\small Khoa Công Nghệ Thông Tin}
    \end{tabular}
  \end{tabular}}
\fancyfoot[L]{\scriptsize\ttfamily \sgufooter}
\fancyfoot[R]{\scriptsize\ttfamily Trang \thepage/\pageref{LastPage}}
\renewcommand{\headrulewidth}{0.3pt}
\renewcommand{\footrulewidth}{0.3pt}

\DefineVerbatimEnvironment{code}{Verbatim}{fontsize=\footnotesize, breaklines=true, frame=lines, framerule=0.3pt, rulecolor=\color{sguline}, numbers=left, numbersep=6pt, xleftmargin=14pt}
\newcommand{\codefile}[3]{\VerbatimInput[fontsize=\footnotesize, breaklines=true, frame=lines, framerule=0.3pt, rulecolor=\color{sguline}, numbers=left, numbersep=6pt, xleftmargin=14pt, firstline=#1, lastline=#2, firstnumber=#1]{#3}}

\newcommand{\fpath}[1]{\nolinkurl{#1}}
\newtcolorbox{ghichu}[1][]{enhanced, breakable, colback=sgugray, colframe=sguline,
  boxrule=0.5pt, arc=2pt, left=8pt, right=8pt, top=5pt, bottom=5pt,
  before skip=8pt, after skip=8pt, #1}

% Trang bìa theo mẫu SGU.  #1 đề tài, #2 loại bài, #3 tháng/năm
\newcommand{\sgutitle}[3]{%
\begin{titlepage}
\begin{center}
TRƯỜNG ĐẠI HỌC SÀI GÒN \\
KHOA CÔNG NGHỆ THÔNG TIN
\end{center}
\vspace{1cm}
\begin{center}\includegraphics[width=3cm]{logoITSGU.png}\end{center}
\vspace{1cm}
\begin{center}
\begin{tabular}{c}
  \multicolumn{1}{l}{\textbf{{\Large CÔNG NGHỆ PHẦN MỀM}}}\\
  ~~\\ \hline \\
  \multicolumn{1}{l}{\textbf{{\Large #2}}}\\ \\
  \textbf{{\Huge #1}}\\ \\ \hline
\end{tabular}
\end{center}
\vspace{2.5cm}
\begin{center}
\begin{tabular}{rrl}
  \hspace{4cm} & GVHD: & Từ Lãng Phiêu\\
  & SV: & \tenLong{} -- \mssvLong\\
  & & \tenTai{} -- \mssvTai\\
\end{tabular}
\end{center}
\vfill
\begin{center}{\footnotesize TP. HỒ CHÍ MINH, #3}\end{center}
\end{titlepage}
\thispagestyle{empty}
\newpage
\tableofcontents
\newpage}

\newcolumntype{L}[1]{>{\raggedright\arraybackslash}p{#1}}

\begin{document}
\sgutitle{BrewLite -- Ứng dụng đặt cà phê}{Assignment 1: Hiện thực nghiệp vụ end-to-end theo Scrum}{THÁNG 11/2026}

% =====================================================================
\section{Giới thiệu}
\subsection{Đề tài và mục tiêu}
BrewLite là ứng dụng đặt đồ uống không dùng tiền mặt. Luồng nghiệp vụ tối thiểu theo đề:
\emph{xem menu $\rightarrow$ chọn size/topping $\rightarrow$ giỏ hàng $\rightarrow$ đặt đơn $\rightarrow$ thanh toán $\rightarrow$ theo dõi đơn}.
Assignment 1 chấm \textbf{việc hiện thực chạy được từ đầu đến cuối} theo quy trình Agile Scrum: 10 task, mỗi task 1 điểm,
công nghệ bắt buộc NestJS (backend) và Next.js (frontend). Các sơ đồ phân tích thiết kế, SRS và kiến trúc thuộc phạm vi
Assignment 2 và không đưa vào báo cáo này.

\subsection{Nhóm và kho mã nguồn}
\begin{table}[H]\centering\small
\begin{tabularx}{\textwidth}{@{}l l X l@{}}
\toprule
\textbf{Thành viên} & \textbf{GitHub} & \textbf{Vai trò} & \textbf{Khu vực mã} \\ \midrule
Nguyễn Thành Long & \texttt{Longcodepro} & Backend, DevOps, Scrum Master & \fpath{backend/}, \fpath{.github/}, \fpath{docker-compose.yml} \\
Trần Tấn Tài & \texttt{TaiTran263} & Frontend, UI/UX & \fpath{frontend/} \\ \bottomrule
\end{tabularx}
\end{table}
Kho mã: \url{https://github.com/Longcodepro/Project_SE_Assignment01_HK1_2026_2027}. Mỗi bài tập có một kho riêng;
kho của Assignment 1 dùng để \emph{kiểm chứng phân công}: mỗi chức năng một nhánh, một Pull Request, một issue,
nên có thể truy từ từng task về người thực hiện.

\subsection{Công nghệ}
\begin{table}[H]\centering\small
\begin{tabularx}{\textwidth}{@{}l X@{}}
\toprule \textbf{Thành phần} & \textbf{Công nghệ} \\ \midrule
Frontend & Next.js 16 (App Router), TypeScript, Tailwind CSS, Zustand (giỏ hàng, phiên đăng nhập), Axios \\
Backend & NestJS 12, TypeScript (ESM), class-validator, JWT, \texttt{bcryptjs}, Swagger tại \texttt{/docs} \\
Cơ sở dữ liệu & PostgreSQL 16, Prisma 7 (7 migration đến 27/10, thêm 3 migration ở mục~\ref{sec:t10}) \\
Thanh toán & Cổng thanh toán mô phỏng (Ví / Thẻ) viết trong \fpath{payments/mock-gateway.ts} \\
Kiểm thử & Vitest + Supertest (unit), test tích hợp trên PostgreSQL thật, Postman/newman \\
DevOps & Docker Compose (3 dịch vụ), GitHub Actions CI/CD, ảnh Docker đẩy lên GHCR \\ \bottomrule
\end{tabularx}
\end{table}

% =====================================================================
\section{Quy trình Scrum}
\subsection{Kế hoạch sprint}
Nhóm chia 3 sprint theo gợi ý của đề. Dữ liệu dưới đây lấy từ issue và milestone trên GitHub
(issue \#3 đến \#19); các issue đều đã đóng.

\begin{table}[H]\centering\small
\begin{tabularx}{\textwidth}{@{}l X l@{}}
\toprule \textbf{Sprint} & \textbf{Mục tiêu và issue} & \textbf{Hoàn thành} \\ \midrule
Sprint 1 & Nền tảng và menu: \#3 CI, \#4 API sản phẩm, \#5 trang Menu & 23/10 -- 25/10/2026 \\
Sprint 2 & Giỏ hàng, đặt đơn, đăng nhập: \#6 chi tiết sản phẩm, \#7 giỏ hàng, \#8 API tạo đơn, \#9 JWT & 26/10/2026 \\
Sprint 3 & Thanh toán, bàn giao, nghiệp vụ nâng cao: \#10 thanh toán, \#11 màn thanh toán, \#12 lịch sử đơn + docker, \#13 màn lịch sử, \#14--\#17 Task 10, \#19 README & 27/10/2026 \\
Ngoài sprint & \#18 CD (đóng gói ảnh Docker lên GHCR, cấu hình Cloudflare Tunnel) & 27/10/2026 \\ \bottomrule
\end{tabularx}
\caption{Issue theo milestone (nguồn: GitHub Issues)}
\end{table}

\subsection{Quy ước làm việc}
\begin{itemize}
  \item Mỗi chức năng một nhánh \texttt{<tên>/task-N-mô-tả}; không đẩy trực tiếp lên \texttt{main}.
  \item Commit dạng \texttt{<loại>(task-N): <mô tả>} bằng tiếng Việt, mỗi commit một mục đích.
  \item Mỗi Pull Request gắn một issue (\texttt{Closes}/\texttt{Refs \#N}); CI phải xanh mới merge.
  \item Trước khi commit backend: \texttt{npm run lint \&\& npm run build \&\& npm test}.
\end{itemize}

% =====================================================================
\section{Hiện thực 10 task}
Bảng sau đối chiếu từng task của đề với hiện thực. Cột PR là Pull Request đã merge vào \texttt{main}.

\begin{table}[H]\centering\footnotesize
\begin{tabularx}{\textwidth}{@{}c L{3.3cm} X l l@{}}
\toprule \textbf{\#} & \textbf{Task của đề} & \textbf{Hiện thực} & \textbf{Người} & \textbf{PR} \\ \midrule
1 & Khởi tạo dự án & Hai thư mục \texttt{frontend}, \texttt{backend}, Docker, README, \texttt{.env.example} & Tài, Long & \#1, \#2 \\
2 & API danh sách sản phẩm & \texttt{GET /products} (Prisma, 6 món mẫu), Swagger, Postman & Long & \#24, \#28 \\
3 & Trang Menu & Lưới sản phẩm; trạng thái loading, rỗng, lỗi + thử lại & Tài & \#22 \\
4 & Chi tiết, tùy chọn & Chọn size S/M/L, topping; giá cập nhật theo tùy chọn & Tài & \#30 \\
5 & Giỏ hàng & Thêm, sửa, xóa; Zustand lưu \texttt{localStorage}; badge số lượng & Tài & \#34 \\
6 & API tạo đơn & \texttt{POST /orders}, validate, đơn \texttt{PENDING}, tổng tiền do server tính & Long & \#31 \\
7 & Đăng ký, đăng nhập & \texttt{POST /auth/register}, \texttt{/auth/login}, bcrypt, JWT, guard; giao diện đăng nhập & Long, Tài & \#32, \#46 \\
8 & Thanh toán không tiền mặt & \texttt{POST /payments} (Ví/Thẻ); màn thanh toán và xác nhận & Long, Tài & \#37, \#35 \\
9 & Xác nhận, lịch sử, bàn giao & \texttt{GET /orders/me}, màn lịch sử, Docker Compose 3 dịch vụ, README & Long, Tài & \#38, \#36, \#43 \\
10 & Nghiệp vụ backend & Xem mục~\ref{sec:t10} & Long & \#39--\#42, \#56--\#58 \\ \bottomrule
\end{tabularx}
\caption{Đối chiếu task và hiện thực}
\end{table}

\subsection{Quy tắc nghiệp vụ đã chốt}
\begin{itemize}
  \item \textbf{Tổng tiền do server tính}: $\sum \text{quantity}\times(\text{giá} + \text{phụ thu size} + \text{giá topping})$.
        Size S/M/L: $+0$/$+5.000$/$+10.000$đ. Topping: trân châu 8.000, kem cheese 10.000, thạch cà phê 6.000đ
        (\fpath{backend/src/orders/pricing.ts}). Client gửi giá lên bị từ chối (400).
  \item \textbf{Thanh toán mô phỏng}: thẻ có 4 số cuối \texttt{0000} bị từ chối; Ví từ chối khi đơn trên 500.000đ.
  \item Mật khẩu băm \texttt{bcryptjs}; không có giá trị mặc định cho \texttt{JWT\_SECRET}.
\end{itemize}

% =====================================================================
\section{Task 10: nghiệp vụ backend}\label{sec:t10}
Đề yêu cầu bốn ý và ba tiêu chí chấm: (a) chặn chuyển trạng thái sai, (b) gửi hai lần cùng \texttt{Idempotency-Key}
chỉ tạo một đơn, (c) đặt đồng thời không vượt tồn kho. Sau khi soát lại mã nguồn ngày 14/11/2026, nhóm nhận thấy
bản hiện thực ngày 27/10 chưa đủ: máy trạng thái chỉ có 3 trạng thái, chưa có tồn kho và mã giảm giá.
Nhóm đã bổ sung ba Pull Request (\#56, \#57, \#58) và chạy lại toàn bộ kiểm thử.

\begin{table}[H]\centering\small
\begin{tabularx}{\textwidth}{@{}L{3.6cm} X l@{}}
\toprule \textbf{Ý của đề} & \textbf{Hiện thực} & \textbf{PR} \\ \midrule
1. State machine & 7 trạng thái, 8 chuyển hợp lệ, \texttt{assertTransition} trả 409 nếu sai & \#39, \#57 \\
2. Idempotency & Header \texttt{Idempotency-Key} cho \texttt{POST /orders} và \texttt{POST /payments}, ràng buộc duy nhất ở DB & \#40 \\
3. Đồng thời / tồn kho & Khóa lạc quan \texttt{version} cho đơn; cột \texttt{Product.stock}, trừ kho nguyên tử & \#41, \#56 \\
4. Khuyến mãi, điểm thưởng & Mã giảm giá \texttt{BREW10}, \texttt{GIAM15K}; tích 1 điểm / 1.000đ & \#42, \#58 \\ \bottomrule
\end{tabularx}
\end{table}

\subsection{State machine đơn hàng}
\begin{figure}[H]\centering
\begin{tikzpicture}[node distance=14mm and 9mm, every node/.style={font=\scriptsize\ttfamily},
  st/.style={draw, rounded corners=2pt, minimum width=19mm, minimum height=7mm, fill=sgugray},
  fin/.style={st, fill=sguline}, ar/.style={-{Stealth[length=2mm]}, thin}]
  \node[st] (pe) {PENDING};
  \node[st, right=of pe] (pa) {PAID};
  \node[st, right=of pa] (pr) {PREPARING};
  \node[st, right=of pr] (rd) {READY};
  \node[fin, right=of rd] (co) {COMPLETED};
  \node[st, below=of pa] (pf) {PAYMENT\_FAILED};
  \node[fin, below=of pe] (ca) {CANCELLED};
  \draw[ar] (pe) -- node[above]{OK} (pa);
  \draw[ar] (pa) -- (pr);
  \draw[ar] (pr) -- (rd);
  \draw[ar] (rd) -- (co);
  \draw[ar] (pe) -- node[below, sloped]{lỗi} (pf);
  \draw[ar] (pf) -- node[right]{trả lại} (pa);
  \draw[ar] (pe) -- node[left]{hủy} (ca);
  \draw[ar] (pf) -- node[above]{hủy} (ca);
\end{tikzpicture}
\caption{Vòng đời đơn (nguồn: \fpath{backend/src/orders/order-status.ts}). Ba bước PAID $\rightarrow$ COMPLETED do nhân viên thực hiện.}
\end{figure}
\codefile{17}{25}{../../backend/src/orders/order-status.ts}

Chỉ nhân viên (cột \texttt{User.role = STAFF}) mới gọi được \texttt{PATCH /orders/:id/status} và \texttt{GET /orders/queue};
khách gọi nhận 403. Vai trò được đọc từ cơ sở dữ liệu ở mỗi lần gọi để việc hạ quyền có hiệu lực ngay.
Đơn đã \texttt{PAID} không hủy được (hoàn tiền nằm ngoài phạm vi).
Mọi lần đổi trạng thái đều qua \texttt{assertTransition} và khóa lạc quan.

\subsection{Chống gửi trùng và khóa lạc quan}
Khóa duy nhất \texttt{@@unique([userId, idempotencyKey])} (đơn) và \texttt{@@unique([orderId, idempotencyKey])} (thanh toán)
làm cho yêu cầu lặp trả lại kết quả cũ. Hai yêu cầu cùng khóa đến cùng lúc: bên thua bắt lỗi \texttt{P2002} và trả kết quả của bên thắng.
Cột \texttt{Order.version} tăng 1 sau mỗi lần ghi; câu lệnh cập nhật có điều kiện \texttt{where: \{ id, version \}}, không khớp thì 409.

\subsection{Tồn kho khi đặt đồng thời}
\codefile{57}{87}{../../backend/src/orders/orders.service.ts}
Phép trừ \texttt{UPDATE ... WHERE stock >= số ly} là một câu lệnh nguyên tử, nên hai giao dịch đồng thời không thể cùng thấy
``còn hàng''. Hết hàng thì ném 409 và giao dịch hoàn lại các món đã trừ. Hủy đơn hoàn lại kho trong cùng giao dịch.
Migration thêm \texttt{CHECK ("stock" >= 0)} làm chốt chặn cuối ở cơ sở dữ liệu.

\subsection{Khuyến mãi và điểm thưởng}
\texttt{POST /orders} nhận \texttt{promoCode} tùy chọn (không phân biệt hoa thường). \texttt{BREW10} giảm 10\%, tối đa 20.000đ;
\texttt{GIAM15K} giảm 15.000đ cho đơn từ 50.000đ. Mã sai hoặc đơn chưa đủ giá trị trả 400 và không trừ kho.
\texttt{total} là số phải trả sau giảm nên thanh toán và điểm thưởng tự tính trên số đã giảm. Điểm thưởng: 1 điểm / 1.000đ,
làm tròn xuống, cộng trong cùng giao dịch với việc chuyển đơn sang \texttt{PAID}; thanh toán lỗi hoặc gửi lại cùng khóa thì không cộng thêm.

% =====================================================================
\section{API}
\begin{table}[H]\centering\footnotesize
\begin{tabularx}{\textwidth}{@{}l l X l@{}}
\toprule \textbf{Phương thức} & \textbf{Đường dẫn} & \textbf{Mô tả} & \textbf{Quyền} \\ \midrule
GET & \texttt{/}, \texttt{/docs} & Kiểm tra backend; Swagger & công khai \\
GET & \texttt{/products} & Menu, kèm \texttt{stock} & công khai \\
POST & \texttt{/auth/register}, \texttt{/auth/login} & Đăng ký, đăng nhập (trả \texttt{accessToken}) & công khai \\
POST & \texttt{/orders} & Đặt đơn; \texttt{Idempotency-Key}, \texttt{promoCode} tùy chọn & khách \\
GET & \texttt{/orders/me} & Lịch sử đơn, mới nhất trước & khách \\
POST & \texttt{/orders/:id/cancel} & Hủy đơn chưa thanh toán, hoàn kho & khách \\
POST & \texttt{/payments} & Thanh toán Ví/Thẻ; \texttt{Idempotency-Key} tùy chọn & khách \\
GET & \texttt{/orders/queue} & Hàng đợi pha chế & nhân viên \\
PATCH & \texttt{/orders/:id/status} & \texttt{PREPARING} $\rightarrow$ \texttt{READY} $\rightarrow$ \texttt{COMPLETED} & nhân viên \\ \bottomrule
\end{tabularx}
\end{table}

% =====================================================================
\section{Kiểm thử và kết quả}
Số liệu dưới đây lấy từ lần chạy ngày 14/11/2026 trên nhánh \texttt{main} sau khi merge PR \#58.

\begin{table}[H]\centering\small
\begin{tabularx}{\textwidth}{@{}l X l@{}}
\toprule \textbf{Hạng mục} & \textbf{Nội dung} & \textbf{Kết quả} \\ \midrule
Lint, build backend & \texttt{oxlint}, \texttt{nest build} & không lỗi \\
Unit test & Vitest, 9 file (auth, orders, payments, products, swagger, state machine, mã giảm giá, loyalty, health) & 93/93 qua \\
Test trên PostgreSQL thật & 3 test tồn kho + 1 test health (\texttt{npm run test:e2e}) & 4/4 qua \\
Postman/newman & 4 request, 12 assertion (health, products, 404, Swagger) & 12/12 qua \\
Độ phủ (v8), toàn backend & dòng / nhánh / hàm & 83,2\% / 75,4\% / 88,4\% \\
Độ phủ, \texttt{orders}, \texttt{payments} & dòng & 97,1\%, 96,2\% \\
CI GitHub Actions & job \texttt{backend}, \texttt{backend-db}, \texttt{frontend}, \texttt{docker} & xanh ở PR \#56, \#57, \#58 \\
CD & đóng gói ảnh Docker backend lên GHCR khi merge vào \texttt{main} & xanh \\ \bottomrule
\end{tabularx}
\end{table}

\subsection{Ba tiêu chí chấm của Task 10}
\begin{enumerate}[label=(\alph*)]
  \item \textbf{Chặn chuyển trạng thái sai}: \fpath{order-status.spec.ts} duyệt đủ $7\times7=49$ cặp trạng thái, chỉ 8 cặp hợp lệ;
        8 trường hợp cấm cụ thể (ví dụ \texttt{PAID}$\rightarrow$\texttt{CANCELLED}, \texttt{PAID}$\rightarrow$\texttt{READY}) trả 409.
  \item \textbf{Gửi hai lần cùng \texttt{Idempotency-Key} chỉ tạo một đơn}: các test \texttt{POST /orders với Idempotency-Key}
        (gửi lại, hai yêu cầu cùng lúc, sai định dạng) và test thanh toán gửi lại không cộng điểm lần nữa.
  \item \textbf{Đặt đồng thời không vượt tồn kho}: \fpath{backend/test/ton-kho.e2e-spec.ts} chạy trên PostgreSQL thật:
        20 khách cùng đặt 1 ly khi kho còn 5 $\Rightarrow$ đúng 5 đơn thành công, 15 yêu cầu nhận 409, kho về 0,
        tổng số ly đã bán bằng 5. Một test khác chứng minh đơn có một món thiếu hàng thì không trừ kho món nào,
        và cơ sở dữ liệu từ chối \texttt{stock = -1}. CI chạy các test này ở job \texttt{backend-db}.
\end{enumerate}

\subsection{Chạy thử luồng thật}
Backend chạy thật trên PostgreSQL tạm (cổng riêng), gọi bằng \texttt{curl}. Kết quả rút gọn:
\begin{Verbatim}[fontsize=\footnotesize, frame=lines, framerule=0.3pt, rulecolor=\color{sguline}]
Đặt 2 ly M + mã brew10   -> 201 PENDING, total 61.200, discount 6.800, promoCode BREW10
Đặt 2 ly khi kho còn 1    -> 409 "Sản phẩm 1 không đủ hàng cho 2 ly"
Mã giảm giá sai           -> 400 "Mã giảm giá SAI không hợp lệ"
Thanh toán ví             -> 201 PAID, pointsEarned 61
Khách gọi API nhân viên   -> 403
Nhân viên PAID->COMPLETED -> 409 "Không thể chuyển đơn từ PAID sang COMPLETED"
Nhân viên PREPARING, READY, COMPLETED -> 200 lần lượt
Hủy đơn đã COMPLETED      -> 409
Đặt 1 ly rồi hủy          -> kho 0 -> 1 (hoàn kho)
Thanh toán đơn đã hủy     -> 409 "Đơn đang ở trạng thái CANCELLED, không thể thanh toán"
\end{Verbatim}

% =====================================================================
\section{Phân công và bằng chứng đóng góp}
Số liệu lấy từ \texttt{git} và GitHub ngày 14/11/2026, tính đến trước 3 Pull Request bổ sung \#56--\#58 (đều của Long).

\begin{table}[H]\centering\small
\begin{tabularx}{\textwidth}{@{}l X r r@{}}
\toprule \textbf{Nội dung} & \textbf{Cách tính} & \textbf{Long} & \textbf{Tài} \\ \midrule
Pull Request đã merge & theo tiền tố nhánh \texttt{long/}, \texttt{tai/}; PR khởi tạo \#1, \#2 tính theo người đặt tên nhánh & 26 & 10 \\
Issue được giao & \#3 đến \#19 (17 issue); \#19 giao cho cả hai & 11 & 5 \\
Commit không tính merge & \texttt{git shortlog -sn --no-merges}, theo tác giả git & 43 & 8 \\
Dòng mã thêm / xóa & \texttt{git log --numstat} trong \fpath{backend/}, \fpath{frontend/} & +18.600 / $-$886 & +11.686 / $-$324 \\ \bottomrule
\end{tabularx}
\caption{Bằng chứng đóng góp theo từng thành viên}
\end{table}

\textbf{Long} phụ trách toàn bộ backend: khởi tạo NestJS, API sản phẩm, đơn hàng, xác thực JWT, thanh toán, lịch sử đơn,
bốn phần của Task 10 cùng phần bổ sung tồn kho, mã giảm giá, vòng đời đơn; CI/CD, Docker Compose, README, Swagger/Postman.
\textbf{Tài} phụ trách frontend: khởi tạo Next.js, trang Menu, chi tiết sản phẩm, giỏ hàng, màn thanh toán và xác nhận, lịch sử đơn.

\begin{ghichu}
\textbf{Lưu ý để đọc số liệu đúng.} Một số Pull Request nhánh \texttt{tai/} (\#46 giao diện đăng nhập, \#48 responsive di động,
\#52 cấu hình URL backend động) có commit mang tài khoản git của Long, nên số commit theo tác giả của Tài thấp hơn khối lượng
frontend thực tế. Vì vậy báo cáo trình bày nhiều cách tính (theo nhánh, theo issue, theo tác giả) và không quy về một tỷ lệ phần trăm.
\end{ghichu}

% =====================================================================
\section{Hạn chế và hướng phát triển}
\begin{itemize}
  \item \textbf{Frontend chưa dùng các API bổ sung ngày 14/11}: nút hủy đơn, ô nhập mã giảm giá, màn hình hàng đợi cho nhân viên,
        hiển thị tồn kho. Các API này đã chạy và có test ở backend; phần giao diện thuộc phạm vi của Tài.
  \item Lịch sử đơn của frontend khi lỗi mạng tự chuyển sang dữ liệu mẫu, chưa phân biệt rõ với dữ liệu thật.
  \item Mã giảm giá cố định trong mã nguồn (\fpath{promo.ts}), chưa có bảng, hạn dùng hay số lượt dùng.
  \item Tồn kho tính theo ly sản phẩm, chưa theo nguyên liệu; hoàn tiền cho đơn đã thanh toán nằm ngoài phạm vi.
  \item Vai trò nhân viên cấp bằng cách sửa cột \texttt{role} trong cơ sở dữ liệu, chưa có giao diện quản trị.
\end{itemize}
Các nội dung phân tích thiết kế (SRS, use-case, activity/sequence/class, kiến trúc, ADR, triển khai production) được làm ở Assignment 2.

% =====================================================================
\begin{thebibliography}{9}
\bibitem{nest} NestJS Documentation, \url{https://docs.nestjs.com}.
\bibitem{prisma} Prisma ORM Documentation, \url{https://www.prisma.io/docs}.
\bibitem{next} Next.js Documentation, \url{https://nextjs.org/docs}.
\bibitem{idem} IETF, \emph{The Idempotency-Key HTTP Header Field} (bản thảo), \url{https://datatracker.ietf.org/doc/draft-ietf-httpapi-idempotency-key-header/}.
\bibitem{de} Đề bài lớn \emph{BrewLite -- Công nghệ phần mềm}, Trường Đại học Sài Gòn, 2026.
\end{thebibliography}
\end{document}
